\chapter{Results, Evidence, and Safe Interpretation}

\section{Evidence rules}
The final evidence package is the sole numerical authority. Feasibility is an
unconditional rate over all trials at the stated point. Other performance
statistics are conditional on physical feasibility unless explicitly described
as paired or common-feasible. These rules avoid silently discarding hard cases.

\section{Core results}
\begin{figure}[H]
\centering
\includegraphics[width=0.93\linewidth]{fig3_cluster_size_tradeoff.png}
\caption{Cluster-size trade-off: feasibility, conditional total transmit power,
and runtime. The lowest observed mean power in the M6 study occurs at an
intermediate sensing-cluster size.}
\end{figure}

Across 30 common M6 scenarios, the proposed method is physically feasible at
each tested $N_{\rm req}\in\{2,3,4,5,6\}$. Mean power is 37.65, 30.83, 29.50,
29.74, and 30.94 mW, respectively. The nonmonotone relation is meaningful:
small clusters lack geometric FIM diversity, while large clusters incur
dedicated-waveform resource and interference costs.

\begin{figure}[H]
\centering
\includegraphics[width=0.93\linewidth]{fig10_method_comparison_vs_nreq.png}
\caption{Association/recovery comparison. FIM-aware selection is a strong
topology baseline; full recovery performs the final coupled robust covariance
optimization and validation.}
\end{figure}

At $N_{\rm req}=3$, mean power is 30.83 mW for the proposed method, 30.93 mW
for FIM-greedy, and 39.29 mW for nearest-AP. The close FIM-greedy result shows
that geometry strongly selects the sensing topology; it does not eliminate the
need for robust beamforming, fixed-topology re-optimization, and certification.

\section{Robustness and QoS sensitivity}
\begin{figure}[H]
\centering
\includegraphics[width=0.86\linewidth]{fig9_csi_robustness.png}
\caption{Sampled CSI robustness. The Monte Carlo result complements, but does
not replace, the worst-case S-procedure formulation.}
\end{figure}

The robust design has zero sampled system outage over the reported uncertainty
radii; the nominal-CSI design has mean outage around 89--92\%. Fig.~8 should be
called a \emph{QoS-requirement sensitivity study}: stricter communication and
tracking targets require more power over the tested grid. It is neither a
global communication--sensing Pareto frontier nor the boundary of the full
feasible region.

\section{Scalability and limitations}
The M12, $K=6$, $P=3$ fixed-cardinality study reports proposed recovery feasible
in 8/8 cases at $N_{\rm req}=3$. Its mean end-to-end runtime is 496.3 seconds
per scenario. This is evidence of larger-system applicability and a clear
motivation for future acceleration, not a real-time claim.

\sourcebox{Detailed bilingual figure explanations and all precise values are in
\texttt{RESULTS\_AND\_FIGURES\_BILINGUAL.md}.}
